\bookchapter{Глава 15: Наука о батареях и оптимизация автономной работы}

\chapterintro{Батарея — это «энергетическое сердце» вашего телефона. Эта глава поможет вам понять характер батареи, освоить научные методы зарядки и продлить время работы, чтобы навсегда избавиться от страха разрядки.}

\booksection{I. Как работает батарея? Циклы и состояние здоровья}

Аккумулятор телефона (Li-ion или Li-Po) похож на энергетическую губку: зарядка — это впитывание воды (накопление энергии), разрядка — выжимание (расход энергии). Важно запомнить несколько особенностей:

\begin{itemize}
  \item \textbf{Заряжайте когда удобно}: в отличие от старых батарей, у литиевых нет эффекта памяти, не обязательно разряжать до нуля.
  \item \textbf{Боится жары}: высокая температура — главный враг батареи.
  \item \textbf{Стареет}: со временем максимальная ёмкость снижается, это нормальный процесс.
\end{itemize}

\booksubsection{Ключевые показатели: количество циклов и здоровье батареи}

\begin{itemize}
  \item Один цикл — это не одна зарядка. Например, сегодня вы израсходовали 60\% и зарядили до 100\%, завтра — 40\% и зарядили. Суммарный расход 100\% считается одним циклом.
  \item Здоровье батареи — это текущая максимальная ёмкость в процентах от новой. При падении ниже 80\% вы заметно почувствуете сокращение времени работы.
  \item \textbf{Как проверить?}iPhone: Настройки > Аккумулятор > Состояние аккумулятора и зарядка; Android: обычно в настройках найдите «Состояние батареи».
\end{itemize}

\booksection{II. Мифы о зарядке? Быстрая, беспроводная и золотое правило}

\begin{itemize}
  \item \textbf{Вредит ли быстрая зарядка батарее?}При использовании оригинальных или сертифицированных аксессуаров вред минимален. Телефон и зарядное устройство имеют термоконтроль — при перегреве мощность автоматически снижается.
  \item \textbf{Беспроводная зарядка}: удобна, но менее эффективна, сильнее нагревается. Подходит для ночной зарядки.
  \item \textbf{Золотое правило сохранения батареи}: старайтесь держать заряд в пределах 20\%–80\%, не допускайте длительного нахождения на 0\% или 100\%.
\end{itemize}

\begin{tcbinfo}{Функция iPhone «Ограничение 80\%» — золотое правило как системная настройка}
В iPhone серии 15/16 (iOS 17 и выше) правило «20\%–80\%» реализовано как системная функция. После включения зарядка автоматически останавливается на 80\%, значительно замедляя старение батареи. Путь настройки:

\texttt{Настройки > Аккумулятор > Состояние аккумулятора и зарядка > Оптимизация зарядки > Лимит 80\%}

Значение трёх вариантов: (1)\textbf{Без ограничений}— зарядка до 100\%; (2)\textbf{Оптимизированная зарядка}— изучает ваш режим, откладывает достижение 100\% на последний этап под утро; (3)\textbf{Лимит 80\%}— останавливается на 80\%, максимально бережёт батарею. Тем, у кого есть постоянный доступ к зарядке дома или в офисе, настоятельно рекомендуется выбрать лимит 80\%; тем, кто часто в разъездах и нуждается в полном заряде, подойдёт оптимизированная зарядка.

\end{tcbinfo}

\begin{tcbinfo}{Зарядные устройства GaN — идеальный спутник в командировках}
Вы, возможно, заметили, что современные блоки быстрой зарядки становятся всё меньше. За это стоит благодарить\textbf{GaN (нитрид галлия)}— новый полупроводниковый материал. По сравнению с традиционными кремниевыми зарядными устройствами, GaN имеет несколько явных преимуществ:

\begin{itemize}
  \item \textbf{Компактность}: при той же мощности GaN-зарядка может быть в 2–3 раза меньше кремниевой. Блок GaN на 65 Вт меньше старого 5-ваттного блока Apple.
  \item \textbf{Высокая эффективность, низкий нагрев}: потери на переключение у GaN намного ниже, чем у кремния, КПД преобразования энергии достигает 95\% и выше, зарядное устройство не перегревается.
  \item \textbf{Идеально для поездок}: одно многопортовое GaN-устройство на 65 или 100 Вт может одновременно заряжать телефон, планшет и ноутбук — гораздо легче, чем возить три «кирпича».
  \item \textbf{Советы по выбору}: выбирайте GaN-зарядку с поддержкой протокола PD (от 65 Вт для ноутбука), подойдут известные бренды: Baseus, Anker, Nubia, Xiaomi и другие. Проверьте совместимость кабеля (для высокой мощности нужен кабель Type-C с чипом E-Marker).
\end{itemize}

\end{tcbinfo}

\booksection{III. Полный разбор протоколов быстрой зарядки: почему ваше зарядное устройство заряжает медленно?}

Вы наверняка сталкивались с ситуацией: телефон поддерживает быструю зарядку, но с другим зарядным устройством скорость резко падает. Это потому, что быстрая зарядка — это не просто «большая мощность», а необходимость, чтобы зарядное устройство, кабель и телефон «говорили на одном языке» — это и есть протокол быстрой зарядки.

\booksubsection{Основной принцип быстрой зарядки}

Мощность зарядки = напряжение × ток (P = U × I). Чтобы увеличить мощность, нужно либо повысить напряжение, либо ток, либо и то и другое. Разные протоколы выбрали разные технические пути.

Процесс быстрой зарядки:

\begin{enumerate}
  \item \textbf{Этап рукопожатия}: телефон и зарядное устройство обмениваются информацией о поддерживаемых протоколах и мощности через кабель.
  \item \textbf{Этап согласования}: обе стороны определяют максимальный общий поддерживаемый уровень мощности.
  \item \textbf{Этап зарядки}: зарядное устройство выдаёт согласованные напряжение и ток, телефон управляет питанием.
  \item \textbf{Этап регулировки}: по мере увеличения заряда телефон постепенно снижает мощность для защиты батареи.
\end{enumerate}

\booksubsection{Подробный обзор основных протоколов быстрой зарядки}

\textbf{USB PD (Power Delivery)}— король универсальных протоколов

\begin{itemize}
  \item \textbf{Особенности}: официальный стандарт организации USB-IF, самый широкий охват. PD 3.0 поддерживает до 100 Вт (20 В/5 А), PD 3.1 расширен до 240 Вт (48 В/5 А).
  \item \textbf{Уровни напряжения}: 5 В, 9 В, 15 В, 20 В, PD 3.1 также поддерживает 28 В, 36 В, 48 В.
  \item \textbf{Совместимость}: вся линейка iPhone, Google Pixel, некоторые флагманы Android, ноутбуки, игровые консоли Switch поддерживают протокол PD.
  \item \textbf{Преимущества}: высокая универсальность, одно зарядное устройство для разных гаджетов; поддержка PPS (программируемый источник питания) для точной регулировки напряжения и тока.
  \item \textbf{Недостатки}: требуется интерфейс Type-C и специальный чип, стоимость выше.
\end{itemize}

\textbf{Qualcomm QC (Quick Charge)}— ветеран мира Android

\begin{itemize}
  \item \textbf{QC 3.0}: регулируемое напряжение от 3,6 до 20 В с шагом 0,2 В, до 18 Вт.
  \item \textbf{QC 4.0/4+}: совместим с USB PD, до 27 Вт, добавлена технология двухканальной зарядки.
  \item \textbf{QC 5.0}: до 100 Вт+, 0–50\% за 5 минут, значительно улучшенный теплоотвод.
  \item \textbf{Совместимость}: большинство телефонов на процессорах Qualcomm Snapdragon поддерживают, но бренды обычно отдают приоритет собственным протоколам.
\end{itemize}

\textbf{OPPO VOOC / OnePlus WARP / realme Dart}— представители низкого напряжения и большого тока

\begin{itemize}
  \item \textbf{Технический путь}: низкое напряжение, большой ток (от 5 В/4 А), меньше нагрев, выше эффективность.
  \item \textbf{VOOC 1.0}: 20 Вт (5 В/4 А), выпущен в 2014 году, положил начало эре быстрой зарядки.
  \item \textbf{SuperVOOC}: 50–65 Вт, двухэлементная последовательная конструкция.
  \item \textbf{SuperVOOC 2.0}: 65 Вт, полная зарядка 4000 м А·ч за 30 минут.
  \item \textbf{150 Вт / 240 Вт Super Flash Charge}: самая быстрая серийная технология быстрой зарядки на данный момент.
  \item \textbf{Особенности}: требуется специальное зарядное устройство и кабель (с дополнительными контактами), несовместим с другими протоколами.
\end{itemize}

\textbf{Huawei SuperCharge}— эволюция высоковольтной быстрой зарядки

\begin{itemize}
  \item \textbf{FCP (Fast Charge Protocol)}: 9 В/2 А, 18 Вт, совместим с QC 2.0.
  \item \textbf{SCP (Super Charge Protocol)}: поддержка более высокой мощности, от 22,5 до 66 Вт и до 100 Вт.
  \item \textbf{Технические особенности}: технология charge pump преобразует высокое напряжение в низкое с большим током для батареи, КПД до 98\%.
  \item \textbf{Совместимость}: только для устройств Huawei/Honor, некоторые зарядные устройства поддерживают PD.
\end{itemize}

\textbf{Xiaomi Surge Charge}— восхождение собственного протокола

\begin{itemize}
  \item \textbf{Чип Surge P1}: собственный чип управления быстрой зарядкой Xiaomi, поддержка проводной зарядки 120 Вт.
  \item \textbf{Чип Surge G1}: чип управления батареей, в паре с P1 обеспечивает более эффективную быструю зарядку.
  \item \textbf{Технические особенности}: двухэлементная последовательная схема + charge pump, поддержка совместимости с несколькими протоколами.
\end{itemize}

\textbf{Быстрая зарядка Apple}— консервативная, но стабильная

\begin{itemize}
  \item \textbf{Протокол}: только USB PD.
  \item \textbf{Мощность}: iPhone 8 и новее поддерживают быструю зарядку, до примерно 27 Вт (iPhone 15 Pro Max может достигать 27 Вт+).
  \item \textbf{Особенности}: относительно консервативная мощность, но отличная оптимизация кривой зарядки и контроль нагрева.
  \item \textbf{Требования к аксессуарам}: требуется кабель USB-C to Lightning (или USB-C to USB-C) и зарядное устройство PD.
\end{itemize}

\booksubsection{Сравнительная таблица протоколов быстрой зарядки}

\begin{center}
\begin{tabularx}{\linewidth}{|X|X|X|X|}
\hline
\textbf{Протокол} & \textbf{Макс. мощность} & \textbf{Технология} & \textbf{Совместимость} \tabularnewline
\hline
USB PD 3.1 & 240 Вт & Высокое напряжение / несколько уровней & Наиболее широкий \tabularnewline
QC 5.0 & 100 Вт+ & Регулируемое высокое напряжение & Модели на Qualcomm \tabularnewline
VOOC/Dart & 240 Вт & Низкое напряжение, большой ток & Только OPPO-серия \tabularnewline
SuperCharge & 100 Вт & Charge Pump & Huawei/Honor \tabularnewline
Быстрая зарядка Apple & \textasciitilde{}27W & Протокол PD & iPhone/iPad \tabularnewline
\hline
\end{tabularx}
\end{center}

\booksubsection{Протоколы беспроводной зарядки}

У беспроводной зарядки тоже есть свои стандарты:

\begin{itemize}
  \item \textbf{Стандарт Qi}: разработан WPC (Wireless Power Consortium), самый распространённый стандарт. Базовый Qi поддерживает 5–15 Вт.
  \item \textbf{Стандарт Qi2}: выпущен в 2023 году, включает магнитное позиционирование (аналог Apple MagSafe), поддержка 15 Вт+.
  \item \textbf{Частные протоколы}: у каждого бренда свои решения для высокоскоростной беспроводной зарядки, например, Xiaomi 50 Вт, Huawei 50 Вт — требуют специального зарядного устройства.
\end{itemize}

\booksubsection{Руководство по выбору быстрой зарядки}

\begin{itemize}
  \item \textbf{Смотрите на совместимость протоколов}: при покупке зарядного устройства убедитесь, что оно поддерживает протокол быстрой зарядки вашего телефона. Сторонние зарядные устройства обычно указывают поддерживаемые протоколы.
  \item \textbf{Смотрите на кабель}: для мощной быстрой зарядки нужен специальный кабель. Обычные кабели могут поддерживать только до 60 Вт.
  \item \textbf{Смотрите на сертификацию}: пользователям Apple рекомендуется выбирать кабели с сертификацией MFi; пользователям Android — оригинальные брендовые или проверенные сторонние.
  \item \textbf{Следите за нагревом}: при быстрой зарядке нагрев телефона нормален, но если он обжигает — это ненормально, рекомендуется приостановить зарядку или сменить зарядное устройство.
\end{itemize}

\begin{tcbwarn}{Советы по безопасности быстрой зарядки}
\begin{itemize}
  \item Избегайте играть в тяжёлые игры во время зарядки — это усиливает нагрев.
  \item Не используйте некачественные зарядные устройства и кабели — это может быть опасно.
  \item Снимайте толстый чехол при зарядке для лучшего теплоотвода.
  \item Если зарядное устройство или телефон аномально нагреваются или пахнут — немедленно прекратите использование.
\end{itemize}

\end{tcbwarn}

\booksection{IV. Новые технологии батарей: кремний-углеродный отрицательный электрод (SiC)}

Если вы следите за презентациями флагманских телефонов 2024–2026 годов, термин «кремний-углеродный анод» стал почти синонимом премиальных батарей. Это важное обновление материала отрицательного электрода литиевых батарей, принцип которого несложен:

\begin{itemize}
  \item \textbf{Традиционный графитовый анод}: в телефонах всегда использовался графит в качестве отрицательного электрода — технология зрелая и стабильная, но плотность энергии близка к теоретическому пределу (около 372 м А·ч/г).
  \item \textbf{Добавление кремния}: теоретическая удельная ёмкость кремния примерно в 10 раз выше, чем у графита (около 4200 м А·ч/г). Добавление небольшого количества кремния в графит создаёт «кремний-углеродный композитный анод», значительно повышая плотность энергии всей батареи.
  \item \textbf{Повышение плотности энергии}: при том же объёме ёмкость кремний-углеродных батарей может увеличиться примерно на\textbf{15\%–20\%}. Именно это позволило флагманам 2024–2026 годов вместить батареи 5500–6000 м А·ч в корпус толщиной 7–9 мм.
\end{itemize}

\begin{tcbinfo}{В каких моделях используется?}
\begin{itemize}
  \item \textbf{Honor}: серия Magic первой внедрила SiC-аноды в коммерческие продукты — Magic6 Pro, Magic7, складной Magic V3.
  \item \textbf{Xiaomi}: Xiaomi 14 Ultra, серия Xiaomi 15 оснащены новым поколением SiC-батарей.
  \item \textbf{vivo}: серии X100 Ultra, X200 Pro используют решения с высокой плотностью SiC-анодов.
  \item \textbf{Huawei/OPPO и другие}: флагманы 2025–2026 годов повсеместно внедряют SiC-аноды — из «чёрной технологии» они становятся стандартом для премиум-сегмента.
\end{itemize}

\end{tcbinfo}

\textbf{Практическое влияние на пользователя}: главное преимущество SiC-анодов — «больше ёмкости в том же объёме». Телефон можно сделать тоньше без ущерба для автономности, либо сохранить толщину корпуса, но увеличить батарею с 4500 до 5500 м А·ч. Иными словами, вам больше не нужно платить «толщиной и весом» за долгую работу. Стоит отметить, что SiC-батареи всё ещё требуют оптимизации управления при низких температурах и полной зарядке через чипы быстрой зарядки — это одна из причин, почему производители рекомендуют долгосрочно держать заряд на уровне 80\%.

\booksection{V. Настоящие советы по экономии заряда}

\begin{itemize}
  \item \textbf{Экран (главный потребитель)}: уменьшите яркость, сократите время автоблокировки. На OLED-экранах тёмный режим существенно экономит заряд.
  \item \textbf{Плохой сигнал}: в лифте или подвале телефон отчаянно ищет сеть, расходуя батарею очень быстро.
  \item \textbf{5G}: если сигнал слабый, переключитесь вручную на 4G — автономность заметно улучшится.
  \item \textbf{Службы геолокации}: для приложений, не связанных с навигацией, установите разрешение геолокации «Только при использовании».
  \item \textbf{Ограничение фоновой активности}: iPhone: Настройки > Основные > Фоновое обновление приложений; Android: в настройках батареи приложений.
\end{itemize}

\begin{tcbpractice}{Проверьте: сделайте осмотр вашей батареи}
\begin{enumerate}
  \item Проверьте здоровье: найдите и запишите максимальную ёмкость вашей батареи.
  \item Анализ расхода: посмотрите графики и список приложений в настройках батареи, найдите главных виновников расхода.
  \item Составьте план: ограничьте фоновую активность приложений с аномальным расходом.
\end{enumerate}

\end{tcbpractice}
